% TODO checklist:
% - [x] Inspected src/observability/metrics.py (core Prometheus metrics wiring)
% - [x] Inspected src/observability/agent_metrics.py (agent-specific metrics)
% - [x] Inspected src/observability/rate_limit_metrics.py (rate limiting metrics)
% - [x] Inspected src/observability/middleware.py (HTTP observability middleware)
% - [x] Inspected ops/prometheus/prometheus.yml (Prometheus configuration)
% - [x] Reviewed Q&A.md sections on metrics and observability
% - [x] Reviewed README.md for metrics overview

\section{Metrics}
\label{sec:metrics}

The Cineca Agentic Platform implements a comprehensive metrics collection framework built on Prometheus, the industry-standard open-source monitoring and alerting toolkit. The metrics subsystem is designed to provide actionable insights into system behavior, support Service Level Objectives (SLOs), and enable data-driven capacity planning and incident response.

\subsection{Metrics Architecture Overview}

The platform's metrics architecture follows a layered design that separates metric definition, collection, and exposition:

\begin{enumerate}
    \item \textbf{Metric Definition Layer}: Metrics are defined using the \texttt{prometheus\_client} Python library, organized into logical groups based on the domain they measure (HTTP requests, agent runs, LLM calls, tools, jobs, rate limiting).
    
    \item \textbf{Collection Layer}: Metrics are collected through instrumentation points distributed across the codebase---middleware for HTTP metrics, service methods for business logic metrics, and adapters for external system interactions.
    
    \item \textbf{Exposition Layer}: The \texttt{/metrics} endpoint, mounted on the FastAPI application, exposes all collected metrics in Prometheus text format for scraping.
\end{enumerate}

The metrics system is initialized through the \texttt{setup\_metrics()} function in \texttt{src/observability/metrics.py}, which creates a dedicated Prometheus registry, registers default collectors (process, platform, garbage collection), and mounts the metrics endpoint:

\begin{lstlisting}[language=Python, caption={Metrics initialization in the FastAPI application}]
def setup_metrics(app: FastAPI) -> None:
    """
    Initialize Prometheus metrics for the provided FastAPI app.
    - Creates/attaches app.state.prometheus_registry
    - Registers default collectors (Process/Platform/GC)
    - Creates/attaches app.state.metrics (MetricStore instance)
    - Mounts /metrics endpoint (GET), excluded from OpenAPI schema
    """
    registry = _build_registry()
    _register_default_collectors(registry)
    store = _MetricStore(registry=registry)
    
    router = APIRouter()
    @router.get("/metrics", include_in_schema=False)
    async def metrics_endpoint() -> Response:
        data = generate_latest(registry)
        return Response(content=data, media_type=CONTENT_TYPE_LATEST)
    
    app.include_router(router)
    app.state.prometheus_registry = registry
    app.state.metrics = store
\end{lstlisting}

\subsection{Multiprocess Mode Support}

The metrics system supports Prometheus multiprocess mode, which is essential when running multiple Uvicorn worker processes behind a process manager like Gunicorn. When the \texttt{PROMETHEUS\_MULTIPROC\_DIR} environment variable is set, the system configures a \texttt{MultiProcessCollector} that aggregates metrics from a shared directory:

\begin{lstlisting}[language=Python, caption={Multiprocess registry configuration}]
def _build_registry() -> CollectorRegistry:
    registry = CollectorRegistry()
    multiproc_dir = os.getenv("PROMETHEUS_MULTIPROC_DIR")
    if multiproc_dir:
        multiprocess.MultiProcessCollector(registry)
        log.info("observability.metrics.multiprocess_enabled", dir=multiproc_dir)
    return registry
\end{lstlisting}

This design ensures accurate metric aggregation in production deployments where horizontal scaling is common.

\subsection{Comprehensive Metrics Catalog}
\label{subsec:metrics-catalog}

The platform exposes a rich set of metrics organized into seven distinct categories, each designed to provide visibility into specific aspects of system behavior.

\subsubsection{HTTP Request Metrics}

HTTP request metrics provide the foundation for understanding API usage patterns, latency distributions, and error rates. These metrics are collected by the \texttt{ObservabilityMiddleware}:

\begin{table}[ht]
\centering
\caption{HTTP request metrics exposed by the platform.}
\label{tab:http-metrics}
\small
\begin{tabularx}{\textwidth}{lllX}
\hline
\textbf{Metric Name} & \textbf{Type} & \textbf{Labels} & \textbf{Description} \\
\hline
\texttt{http\_requests\_total} & Counter & method, path, status & Total HTTP requests processed \\
\texttt{http\_request\_duration\_seconds} & Histogram & method, path, status & Request latency distribution \\
\hline
\end{tabularx}
\end{table}

The duration histogram uses carefully chosen buckets optimized for API latencies: 5ms, 10ms, 25ms, 50ms, 100ms, 250ms, 500ms, 1s, 2.5s, 5s, and 10s. The \texttt{path} label uses parameterized route templates (e.g., \texttt{/items/\{item\_id\}}) rather than actual paths to avoid high-cardinality label explosion.

\subsubsection{Agent Run Metrics}

Agent-specific metrics track the execution of multi-step agent workflows, providing insights into orchestration performance and resource consumption:

\begin{table}[ht]
\centering
\caption{Agent orchestration metrics.}
\label{tab:agent-metrics}
\small
\begin{tabularx}{\textwidth}{lllX}
\hline
\textbf{Metric Name} & \textbf{Type} & \textbf{Labels} & \textbf{Description} \\
\hline
\texttt{agent\_runs\_total} & Counter & agent\_type, status, tenant\_id & Total agent runs initiated \\
\texttt{agent\_run\_duration\_seconds} & Histogram & agent\_type, status & End-to-end run duration \\
\texttt{agent\_active\_runs} & Gauge & agent\_type, tenant\_id & Currently active runs \\
\texttt{agent\_phase\_duration\_seconds} & Histogram & agent\_type, phase & Duration per execution phase \\
\texttt{agent\_errors\_total} & Counter & agent\_type, error\_type, phase & Errors by type and phase \\
\texttt{agent\_retries\_total} & Counter & agent\_type, reason & Retry attempts \\
\texttt{agent\_queue\_depth} & Gauge & priority & Execution queue depth \\
\texttt{agent\_concurrency\_limit} & Gauge & tenant\_id & Max concurrent runs allowed \\
\texttt{agent\_concurrency\_throttled\_total} & Counter & tenant\_id & Throttled run requests \\
\hline
\end{tabularx}
\end{table}

The agent metrics are particularly valuable for understanding multi-tenant resource consumption patterns and identifying performance bottlenecks in the orchestration pipeline.

\subsubsection{LLM Call Metrics}

LLM call metrics provide visibility into interactions with language model providers, including latency, token consumption, and error tracking:

\begin{table}[ht]
\centering
\caption{LLM provider interaction metrics.}
\label{tab:llm-metrics}
\small
\begin{tabularx}{\textwidth}{lllX}
\hline
\textbf{Metric Name} & \textbf{Type} & \textbf{Labels} & \textbf{Description} \\
\hline
\texttt{llm\_calls\_total} & Counter & model, provider, status & Total LLM API calls \\
\texttt{llm\_call\_duration\_seconds} & Histogram & model, provider, status & API call latency \\
\texttt{llm\_tokens\_total} & Counter & model, provider, type & Tokens consumed (prompt/completion/total) \\
\texttt{llm\_errors\_total} & Counter & model, provider, error\_type & LLM API errors \\
\texttt{model\_requests\_total} & Counter & route, provider, instance, status\_class & Model runtime calls \\
\texttt{model\_request\_latency\_ms} & Histogram & route, provider, instance, status\_class & Model latency (milliseconds) \\
\hline
\end{tabularx}
\end{table}

These metrics enable cost attribution across tenants, capacity planning for LLM infrastructure, and early detection of provider degradation.

\subsubsection{Tool Invocation Metrics}

MCP tool invocation metrics track the execution of the platform's 34 integrated tools:

\begin{table}[ht]
\centering
\caption{Tool invocation metrics.}
\label{tab:tool-metrics}
\small
\begin{tabularx}{\textwidth}{lllX}
\hline
\textbf{Metric Name} & \textbf{Type} & \textbf{Labels} & \textbf{Description} \\
\hline
\texttt{tools\_invocations\_total} & Counter & tool\_name, status, tenant\_id & Total tool invocations \\
\texttt{tools\_invocation\_duration\_seconds} & Histogram & tool\_name, status & Invocation duration \\
\texttt{tools\_queue\_depth} & Gauge & tool\_name & Pending invocations \\
\texttt{tools\_cache\_operations\_total} & Counter & operation, result & Redis cache operations \\
\texttt{tools\_idempotency\_conflicts\_total} & Counter & tool\_name & Idempotency key conflicts (409) \\
\texttt{agent\_tool\_calls\_total} & Counter & agent\_type, tool\_name, status & Tool calls within agent context \\
\texttt{agent\_tool\_call\_duration\_seconds} & Histogram & tool\_name, status & Tool duration in agent context \\
\hline
\end{tabularx}
\end{table}

\subsubsection{Job Metrics}

Background job metrics track the asynchronous job processing system:

\begin{table}[ht]
\centering
\caption{Background job processing metrics.}
\label{tab:job-metrics}
\small
\begin{tabularx}{\textwidth}{lllX}
\hline
\textbf{Metric Name} & \textbf{Type} & \textbf{Labels} & \textbf{Description} \\
\hline
\texttt{background\_jobs\_total} & Counter & job, status & Jobs executed by outcome \\
\texttt{background\_job\_duration\_seconds} & Histogram & job, status & Job execution duration \\
\texttt{job\_create\_total} & Counter & backend, status & Job creation operations \\
\texttt{job\_create\_duration\_seconds} & Histogram & backend & Job creation latency \\
\texttt{job\_get\_duration\_seconds} & Histogram & backend & Job retrieval latency \\
\texttt{sse\_connections\_active} & Gauge & --- & Active SSE connections \\
\texttt{sse\_gap\_events\_total} & Counter & --- & SSE ring buffer evictions \\
\hline
\end{tabularx}
\end{table}

\subsubsection{Intent Classification Metrics}

Intent classification metrics track the performance and accuracy of the intent classification pipeline:

\begin{table}[ht]
\centering
\caption{Intent classification metrics.}
\label{tab:intent-metrics}
\small
\begin{tabularx}{\textwidth}{lllX}
\hline
\textbf{Metric Name} & \textbf{Type} & \textbf{Labels} & \textbf{Description} \\
\hline
\texttt{intent\_classification\_total} & Counter & mode, source, adjusted & Classifications performed \\
\texttt{intent\_classification\_duration\_seconds} & Histogram & mode, source & Classification latency \\
\texttt{intent\_classification\_confidence} & Histogram & mode, source & Confidence score distribution \\
\texttt{intent\_pattern\_matches\_total} & Counter & pattern\_group & Pattern matches by group \\
\texttt{intent\_llm\_fallback\_total} & Counter & success & LLM fallback attempts \\
\texttt{intent\_rbac\_adjustments\_total} & Counter & original\_mode, adjusted\_mode, role & RBAC intent adjustments \\
\hline
\end{tabularx}
\end{table}

\subsubsection{Rate Limiting Metrics}

Rate limiting metrics provide visibility into quota enforcement and abuse detection:

\begin{table}[ht]
\centering
\caption{Rate limiting metrics.}
\label{tab:rate-limit-metrics}
\small
\begin{tabularx}{\textwidth}{lllX}
\hline
\textbf{Metric Name} & \textbf{Type} & \textbf{Labels} & \textbf{Description} \\
\hline
\texttt{rate\_limit\_requests\_total} & Counter & action, scope, result & Rate limit checks (allowed/blocked) \\
\texttt{rate\_limit\_exceeded\_total} & Counter & action, scope & Rate limit violations \\
\texttt{tenant\_quota\_exceeded\_total} & Counter & action, tenant\_id & Tenant quota violations \\
\texttt{rate\_limit\_usage\_ratio} & Histogram & action, scope & Usage ratio (current/limit) \\
\hline
\end{tabularx}
\end{table}

\subsection{Metrics Recording Patterns}

The platform provides a set of helper functions that encapsulate the complexity of metric recording while ensuring consistent labeling and error handling:

\begin{lstlisting}[language=Python, caption={Example of defensive metric recording}]
def record_tool_invocation(
    tool_name: str,
    status: str,
    duration_seconds: float | None = None,
    tenant_id: str = "default-tenant",
    app: FastAPI | None = None,
) -> None:
    """Update tool invocation counters/histograms if metrics were set up."""
    store = _get_store(app)
    if not store:
        return  # No-op if metrics not configured
    try:
        store.tools_invocations_total.labels(tool_name, status, tenant_id).inc()
        if duration_seconds is not None:
            store.tools_invocation_duration_seconds.labels(
                tool_name, status
            ).observe(max(0.0, float(duration_seconds)))
    except Exception as e:
        log.warning("observability.metrics.record_tool_invocation_failed", 
                    error=str(e))
\end{lstlisting}

This defensive pattern ensures that metric recording failures never impact request processing, while still logging any issues for debugging.

\subsection{Service Information Metric}

The platform exposes a static \texttt{service\_info} gauge that provides version information as metric labels:

\begin{lstlisting}[language=Python, caption={Service information metric}]
self.service_info = Gauge(
    "service_info",
    "Static service information (value is always 1).",
    ["version"],
    registry=registry,
)
self.service_info.labels(version=__version__).set(1)
\end{lstlisting}

This pattern allows operators to query deployment versions across instances and track rolling deployments through Prometheus queries.

\subsection{Prometheus Scraping Configuration}

The platform ships with a production-ready Prometheus configuration that scrapes the FastAPI application's \texttt{/metrics} endpoint:

\begin{lstlisting}[language=yaml, caption={Prometheus scrape configuration (ops/prometheus/prometheus.yml)}]
global:
  scrape_interval: 15s
  scrape_timeout: 10s
  evaluation_interval: 30s
  external_labels:
    service: cineca-agentic-platform
    env: dev

scrape_configs:
  - job_name: "app"
    honor_labels: true
    scheme: http
    metrics_path: /metrics
    static_configs:
      - targets: ["app:8000"]
        labels:
          instance: "app"
          component: "fastapi"
\end{lstlisting}

The configuration is tuned for development deployments with 15-second scrape intervals; production deployments may adjust these intervals based on metric cardinality and storage constraints.

